\documentclass[conference]{IEEEtran}

\usepackage{cite}
\usepackage{amsmath,amssymb,amsfonts}
\usepackage{graphicx}
\usepackage{booktabs}
\usepackage{array}
\usepackage{url}
\usepackage{xcolor}

\title{Notes on Polynomial Regression}

\author{
\IEEEauthorblockN{Luis Mendez}
}

\begin{document}
\maketitle

\section*{Definition}
Polynomial regression is a type of regression analysis in which the relationship between the independent variable \( x \) and the dependent variable \( y \) is modeled as an \( n \)-th degree polynomial. The general form of a polynomial regression model is:
\[y = \beta_0 + \beta_1 x + \beta_2 x^2 + \ldots + \beta_n x^n + \epsilon\]
where \( \beta_0, \beta_1, \ldots, \beta_n \) are the coefficients of the polynomial, and \( \epsilon \) is the error term.
\section*{Difference between linear regression and polynomial regression}
Linear regression models the relationship between variables as a straight line, while polynomial regression models the relationship as a curve of a specified degree.
\section*{When to use polynomial regression}
Polynomial regression is useful when the relationship between the independent and dependent variables is non-linear. It can
capture more complex patterns in the data compared to linear regression.

\end{document}